\section{Limitations and Future Work}\label{sec:limitations}

\subsection{Limitations of the measurement}
\textbf{(a) The reference is not an independent ground truth.} It is fused from Faro depth with the same pipeline
(\SI{1}{cm} voxels) and uses the \textbf{same poses} as every row, so (i) the \code{gt} row is non-zero and measures our
own discretisation, (ii) VIO pose drift is removed from every number---a real phone system incurs additional pose
error this paper does not measure, and (iii) areas Faro did not see are masked from every source equally, so
``completeness'' means ``complete relative to what Faro saw,'' not the whole room.

\textbf{(b) Evaluated at $\sim$4~fps only.} Faro depth exists only for the frames Apple kept, so every source is
evaluated on that set although LiDAR runs at 60~fps and RGB at 30~fps. What a real pipeline gains or loses from denser
frames (smoothing vs.\ accumulated error) is not measured; Exp.~3 answers coverage, and compute-vs-accuracy only partly.

\textbf{(c) Fusion resolution $256\times192$}, chosen to be fair to LiDAR, downsamples the monocular model (which predicts
at $640\times480$) before fusion; DA-v2's sharp-edge advantage does not appear in the 3D numbers. ``Mono at its best''
would require fusing at higher resolution and accepting a disadvantaged LiDAR row.

\textbf{(d) Six rooms (3D) / 26 rooms (2D), one dataset, one device.} Every mesh number comes from the six rooms that
have a reference mesh; Exp.~7 extends the 2D measurement of the main result to 6 $+$ 20 rooms but says nothing about
mesh quality there, and uses every second Faro frame. Every scan is an iPad~Pro in homes Apple selected; offices, empty
rooms and low light are absent, and 42444474 was used during development, so its numbers are in-sample---per-scene
tables show whether the other five confirm.

\textbf{(e) The fine-tuned rows are selected on three validation scenes.} Exp.~6 keeps the checkpoint with the lowest
AbsRel on the 3 validation rooms; that ranking (R3 $>$ R2 $>$ R1) does not reproduce on the six test rooms
(R1 $\geq$ R3 $\geq$ R2 in Chamfer). Three rooms are too few to separate checkpoints that differ by $\sim$0.01 AbsRel,
so the Exp.~6 ordering between teachers should be read as ``no measurable difference'', and a larger validation set is
the first thing to add. The six test rooms were never used for training, validation or checkpoint choice. Moreover R1
and R2 do not differ in the label alone: R2 sees $\sim$2.4$\times$ more distinct views (\num{8656} vs.\ \num{3594}
frames at equal steps, \S\ref{sec:ftsetup}); a controlled teacher comparison needs R2 trained on R1's Faro timestamps,
which we have not run. This does not affect the main question (R3 vs.\ pretrained), where both rows are scored on the
same test frames.

\textbf{(f) The no-fine-tuning baseline is one model, one implementation, one input resolution.} Depth Pro is run through
its HuggingFace implementation on $640\times480$ frames upsampled to $1536^2$ by its own processor, with the true
focal length supplied; we did not test Apple's reference code, higher-resolution input (the dataset's
$1920\times1440$ frames were not downloaded), or the other intrinsics-conditioned models (UniDepth, Metric3D).
The claim we make is about this model on this input, not about the family.

\textbf{(g) One fine-tuning recipe, one architecture, one label preparation.} The LiDAR label ($256\times192$) is
up-sampled with nearest-neighbour to the colour grid ($640\times480$, then $518\times686$) before the loss, rather than
down-sampling the prediction to LiDAR's own grid, and confidence $\geq 1$ (not $=2$) is used; neither choice has been
ablated (R2 landing close to R1, whose Faro labels are high-resolution, suggests the effect is small but does not prove
it). Only DA-v2 Metric-Indoor Large with a frozen encoder was fine-tuned; whether unfreezing the encoder, a different loss or another architecture would widen or close the
Faro-vs-LiDAR teacher gap is untested, as is training on more than 24 rooms.

\textbf{(h) Times are Apple Silicon (MPS) with unoptimised PyTorch}, not representative of on-device inference
(CoreML / NNAPI); use them for \emph{ratios} between models, not absolute values.

\subsection{Limitations of the system}
\textbf{(a) \code{per\_scene} still calibrates on GT}: it is the upper bound of one-off calibration; a real system
using a measured wall or camera height would do worse. \S\ref{sec:drift} shows even this bound is far from the oracle
because scale swings per view---the assumption, not the calibration method, is the problem.

\textbf{(b) \code{sparse\_points} is evaluated on a proxy.} ARKitScenes ships no ARKit feature points, so the aligner sees
200 pixels sampled from the LiDAR frame at high confidence. Real VIO points are triangulated, noisier, and cluster on
texture; the row is therefore an \emph{upper bound} of a VIO-based system. The loader exposes a noise knob
(\code{sparse\_noise}) and point count for sensitivity studies; an evaluation on real ARKit/ARCore point clouds needs
our own captures (\code{dataio/custom.py} reads them).

\textbf{(c) No outlier rejection before fusion}: every posed frame is integrated equally. \S\ref{sec:drift} shows a
minority of frames (close-ups, blank walls) create most of \code{per\_scene}'s error; frame weighting is what Exp.~5's
mechanism provides, applied here only to LiDAR confidence so the main tables measure the depth source alone.

\textbf{(d) Poses come from the dataset}: there is no tracking in the system; an MVP uses ARKit/ARCore poses of similar
quality to the trajectories used here, possibly worse on low-end Android.

\textbf{(e) No GPU fusion}: fusion is $<$\SI{2}{s} per scan on CPU, so it is not a bottleneck; a GPU backend fits behind
the same \code{TSDFFusion} interface if higher resolutions are pursued.

\subsection{Future work (ordered by impact on the MVP decision)}
\begin{enumerate}
\item \textbf{R3 $+$ per-frame scale (\code{sparse\_points} on the fine-tuned model)}: \S\ref{sec:whyft} shows that
almost all of R3's remaining error is a per-frame factor (AbsRel 0.134 $\rightarrow$ 0.051 with the right scale); one
new Exp.~6 row answers how close ``LiDAR teacher $+$ VIO points'' gets to a sensor. Together with a \emph{controlled
teacher comparison} (R2 on R1's Faro timestamps) and a \emph{larger validation set}.
\item \textbf{Training data from more phone cameras}: R3 has only seen the iPad camera; it must be re-measured on
iPhone / Android cameras before claiming it serves devices without LiDAR in general.
\item \textbf{Real VIO points}: capture rooms with the app (ARKit \code{rawFeaturePoints} / ARCore point cloud) plus a
reference (tape measurements or a rented laser scanner) and rerun the \code{sparse\_points} row on real, noisy, textured
points; sweep \code{sparse\_noise} and point count on ARKitScenes meanwhile.
\item \textbf{Metric model $+$ scale-only per frame}: \S\ref{sec:metric} suggests one parameter per frame suffices;
test whether one VIO point per frame is enough.
\item \textbf{Frame selection / weighting} for mono: drop or down-weight frames with median depth $<$\SI{1}{m} or low
disparity variance before fusion, using the per-pixel weight channel already used in Exp.~5.
\item \textbf{Android capture (ARCore depth source)}: \code{depth\_sources/arcore.py} has its registry slot; data must be
collected with a reference.
\item \textbf{Compressed models}: Exp.~4 provides the L / S / MiDaS baseline; a distilled or quantised model plugs in as
\code{models/<x>.py} and reruns Exp.~4.
\item \textbf{Fusion above $256\times192$ for mono}, to measure what downsampling costs DA-v2.
\item \textbf{Semantic layer} (wall / floor / furniture labels on the mesh) for usable-area computation; sits after
post-processing and does not affect the geometry numbers.
\end{enumerate}
